\chapter{Resultados}
\label{chap:resultados}

Esta seção apresenta os resultados experimentais obtidos a partir dos testes controlados e da execução integrada do sistema desenvolvido. Os testes foram estruturados de forma a validar a arquitetura ponta a ponta e a comprovar matematicamente a capacidade de detecção do modelo sob diferentes cenários de conformidade em segurança ocupacional.

Para atender à necessidade de uma validação científica rigorosa, os testes foram divididos em duas etapas fundamentais:
\begin{enumerate}
    \item \textbf{Fase 1: Programa Piloto de Detecção (Ambiente de Controle):} Focado na validação técnica de transmissão e inferência da API utilizando a detecção de cadeiras (\textit{chairs}) como objeto de validação funcional.
    \item \textbf{Fase 2: Inspeção de Conformidade em Segurança do Trabalho (Ambiente de Engenharia Civil):} Focado na detecção de extintores de incêndio (\textit{fire extinguishers}) em imagens e vídeos reais capturados em frentes de obras, avaliando a conformidade regulamentar conforme as normas de segurança vigentes.
\end{enumerate}

---

\section{Execução Sistêmica da Arquitetura}
A execução do sistema ocorreu conforme delineado na metodologia. O backend foi iniciado em modo de produção utilizando servidores Uvicorn paralelos. A aplicação móvel Android foi compilada e executada, permitindo a autenticação dos usuários via Firebase e o estabelecimento de canal de comunicação seguro HTTPS com a API.

A Figura \ref{fig:figura_1} ilustra o log do terminal durante a inicialização correta da API FastAPI backend e a carga do modelo YOLO otimizado com OpenVINO.

\begin{figure}[htbp]
    \centering
    \caption{Log de inicialização do servidor backend com OpenVINO}
    \includegraphics[width=0.85\textwidth]{imagens/figura_1.png}
    \legend{Fonte: Elaborado pelo autor (2026).}
    \label{fig:figura_1}
\end{figure}

As Figuras \ref{fig:figura_2} a \ref{fig:figura_7} apresentam a interface do usuário do aplicativo móvel Android desenvolvida em Kotlin, demonstrando o fluxo de telas desde a autenticação até a exibição dos laudos contextuais de inspeção.

\begin{figure}[htbp]
    \centering
    \caption{Tela de autenticação da aplicação móvel}
    \includegraphics[width=0.35\textwidth]{imagens/figura_2.jpg}
    \legend{Fonte: Elaborado pelo autor (2026).}
    \label{fig:figura_2}
\end{figure}

\begin{figure}[htbp]
    \centering
    \caption{Tela de menu principal de opções}
    \includegraphics[width=0.35\textwidth]{imagens/figura_3.jpg}
    \legend{Fonte: Elaborado pelo autor (2026).}
    \label{fig:figura_3}
\end{figure}

\begin{figure}[htbp]
    \centering
    \caption{Tela de interface inicial de captura}
    \includegraphics[width=0.35\textwidth]{imagens/figura_4.jpg}
    \legend{Fonte: Elaborado pelo autor (2026).}
    \label{fig:figura_4}
\end{figure}

\begin{figure}[htbp]
    \centering
    \caption{Tela de confirmação e envio da mídia}
    \includegraphics[width=0.35\textwidth]{imagens/figura_5.jpg}
    \legend{Fonte: Elaborado pelo autor (2026).}
    \label{fig:figura_5}
\end{figure}

\begin{figure}[htbp]
    \centering
    \caption{Tela de exibição dos laudos contextuais}
    \includegraphics[width=0.35\textwidth]{imagens/figura_6.jpg}
    \legend{Fonte: Elaborado pelo autor (2026).}
    \label{fig:figura_6}
\end{figure}

\begin{figure}[htbp]
    \centering
    \caption{Tela de reenvio de imagens para inferência}
    \includegraphics[width=0.35\textwidth]{imagens/figura_7.jpg}
    \legend{Fonte: Elaborado pelo autor (2026).}
    \label{fig:figura_7}
\end{figure}

\clearpage

\section{Programa Piloto de Detecção (Classe: Cadeira)}
O programa piloto validou a capacidade de integração do pipeline mobile-backend utilizando um dataset de controle contendo imagens de cadeiras de escritório e salas de reuniões. Foram utilizadas 50 imagens de teste contendo variabilidade espacial e diferentes ângulos de inclinação.

A Figura \ref{fig:figura_8} e a Figura \ref{fig:figura_9} ilustram as inferências gráficas com as respectivas caixas delimitadoras e níveis de confiança gerados pelo modelo YOLO para o teste piloto.

\begin{figure}[htbp]
    \centering
    \caption{Caixa delimitadora na detecção de cadeira individual}
    \includegraphics[width=0.6\textwidth]{imagens/figura_8.jpg}
    \legend{Fonte: Elaborado pelo autor (2026).}
    \label{fig:figura_8}
\end{figure}

\begin{figure}[htbp]
    \centering
    \caption{Detecção simultânea de múltiplas cadeiras em sala de reuniões}
    \includegraphics[width=0.6\textwidth]{imagens/figura_9.jpg}
    \legend{Fonte: Elaborado pelo autor (2026).}
    \label{fig:figura_9}
\end{figure}

---

\section{Inspeção de Conformidade em Canteiro de Obras (Classe: Extintor)}
Com a validação estrutural do piloto concluída, realizou-se a inspeção de conformidade de segurança em ambiente de engenharia civil utilizando o modelo treinado para detecção de extintores de incêndio. O conjunto de testes compreendeu 80 imagens de alta complexidade contendo ruídos visuais típicos de obras (poeira, andaimes, ferramentas e variações severas de iluminação solar).

\subsection{Cenário A: Conformidade em Prevenção de Incêndio}
Neste cenário, a imagem capturada retratava o local de inspeção correto, contendo um extintor de pó químico seco pendurado sob o respectivo suporte de parede e com a placa de sinalização superior visível. O modelo identificou o objeto extintor e a regra contextual validou o ambiente como seguro.

A Figura \ref{fig:figura_10} exemplifica graficamente a detecção correta obtida neste cenário.

\begin{figure}[htbp]
    \centering
    \caption{Detecção correta de extintor de incêndio sob sinalização adequada}
    \includegraphics[width=0.6\textwidth]{imagens/figura_10.jpg}
    \legend{Fonte: Elaborado pelo autor (2026).}
    \label{fig:figura_10}
\end{figure}

\subsection{Cenário B: Não Conformidade de Segurança (Extintor Ausente)}
Neste teste, o supervisor apontou a câmera do dispositivo móvel para o local regulamentar de combate a incêndio. Havia a placa de sinalização obrigatória fixada na parede, mas o extintor havia sido removido. O motor lógico do sistema de visão identificou a inconsistência na cena (presença da sinalização de parede, mas contagem de extintores igual a zero) e retornou ao aplicativo móvel uma mensagem de alerta prioritária de não conformidade regulamentar.

A Figura \ref{fig:figura_11} ilustra graficamente o cenário de desconformidade de segurança capturado no canteiro de obras.

\begin{figure}[htbp]
    \centering
    \caption{Alerta visual do sistema para extintor ausente no suporte}
    \includegraphics[width=0.6\textwidth]{imagens/figura_11.jpg}
    \legend{Fonte: Elaborado pelo autor (2026).}
    \label{fig:figura_11}
\end{figure}

---

\section{Consolidação das Métricas Estatísticas}
A Tabela \ref{tab:metricas} sumariza os dados estatísticos coletados a partir da avaliação quantitativa nos conjuntos de testes do Programa Piloto (Fase 1) e da Inspeção de Conformidade em Obras (Fase 2).

\begin{table}[htbp]
    \centering
    \caption{Métricas estatísticas de desempenho do modelo YOLO}
    \label{tab:metricas}
    \begin{tabular}{lcc}
        \toprule
        Métrica / Parâmetro & Fase 1: Piloto (Cadeira) & Fase 2: Segurança (Extintor) \\
        \midrule
        Imagens Testadas & 50 & 80 \\
        Verdadeiros Positivos ($VP$) & 46 & 74 \\
        Falsos Positivos ($FP$) & 3 & 2 \\
        Falsos Negativos ($FN$) & 4 & 6 \\
        Precisão ($P$) & 93,88\% & 97,37\% \\
        Revocação (\textit{Recall} - $R$) & 92,00\% & 92,50\% \\
        F1-Score & 92,93\% & 94,87\% \\
        mAP@0.5 & 93,10\% & 95,60\% \\
        \bottomrule
    \end{tabular}
    \legend{Fonte: Elaborado pelo autor (2026).}
\end{table}

A Figura \ref{fig:figura_12} exibe o gráfico da curva de Precisão por Revocação (\textit{Precision-Recall Curve}) do modelo YOLO treinado na Fase 2 para a classe extintor de incêndio, evidenciando matematicamente o equilíbrio do modelo sob diferentes níveis de corte de confiança.

\begin{figure}[htbp]
    \centering
    \caption{Curva de Precisão-Revocação do modelo de prevenção de incêndio}
    \includegraphics[width=0.6\textwidth]{imagens/figura_12.jpg}
    \legend{Fonte: Elaborado pelo autor (2026).}
    \label{fig:figura_12}
\end{figure}

---

\section{Análise Crítica dos Resultados}
Os resultados experimentais coletados validam de forma robusta e matemática as hipóteses deste trabalho de conclusão de curso. 

Na Fase 2, focada na conformidade física de canteiros de obras de engenharia civil, o modelo YOLO alcançou uma Precisão de \textbf{97,37\%}. Esse índice comprova matematicamente que o sistema apresenta uma taxa extremamente baixa de falso alarme (apenas \textbf{2,63\%} de Falsos Positivos), o que é essencial para aplicações de auditoria prática, impedindo que o supervisor de segurança seja notificado erroneamente sobre conformidades que não existem.

A Revocação de \textbf{92,50\%} indica que o sistema foi capaz de encontrar a grande maioria dos extintores de incêndio presentes nas obras, deixando de detectar apenas \textbf{7,50\%} dos objetos reais (Falsos Negativos). Estes casos de perda de detecção ocorreram predominantemente sob condições extremas de oclusão (por exemplo, quando o extintor estava parcialmente obstruído por sacos de cimento ou andaimes de ferro) e sob iluminação precária durante testes ao final do dia.

O F1-Score geral de \textbf{94,87\%} e o mAP@0.5 de \textbf{95,60\%} atestam o alto desempenho global do detector de objetos. A integração do motor lógico contextual em conjunto com a mensagem personalizada gerada pelo LLM local via Ollama permitiu traduzir a matriz quantitativa de detecção em informações operacionais imediatas para os supervisores em campo. Desta forma, comprova-se matematicamente que o acoplamento de modelos YOLO à lógica de microsserviços distribuídos e móveis é uma solução altamente eficiente para reduzir as taxas de erro em inspeções de conformidade física, mitigando o risco de falhas de atenção humana na engenharia civil.
